\documentclass{article}

\textwidth=6.5in
\textheight=8.5in
\voffset=-54pt
\oddsidemargin=0in
\evensidemargin=0in
\setlength{\parskip}{8pt}
\setlength{\parindent}{0pt}

\usepackage{amsmath, amssymb}
\usepackage{ifthen}

\usepackage[inline]{enumitem}
\setlist{topsep=0pt, parsep=8pt}

\usepackage[rgb,table]{xcolor}\convertcolorsUfalse
\usepackage{graphicx}

% change hyperref options
\PassOptionsToPackage{hyphens}{url}
\usepackage[bookmarks,colorlinks,linkcolor=black,citecolor=blue,pdfstartview=FitH,urlcolor=blue]{hyperref}
\hypersetup{pdfpagemode=UseNone}
\usepackage{bookmark}

\usepackage{graphicx}
\usepackage{multicol}
\usepackage{framed}
\usepackage{array}

\usepackage{icomma}

\usepackage{pifont}

\usepackage{soul}

\usepackage{pgfplots}
\pgfplotsset{compat=1.18}  
\pgfplotscreateplotcyclelist{mycolor}{black, blue, red, green!70!black, magenta, purple, cyan, orange }  
\pgfplotsset{
  disabledatascaling,
  cycle list=tikzcolor, 
  axis lines=middle,
  every axis plot/.style={no marks, thick},
  every axis/.style={
    enlarge x limits=false,
    enlarge y limits=auto,
    xlabel style={at={(current axis.right of origin)},anchor=west},
    ylabel style={at={(current axis.above origin)},anchor=south},
    % extra x and y ticks for asymptotes
    extra x tick style={grid=major, grid style={dashed,black}},
    extra y tick style={grid=major, grid style={dashed,black}},
    /pgf/number format/1000 sep={},
  },    
  tick label style = {font=\footnotesize, fill=\currentbgcolor, inner sep=0pt},    
  xticklabel shift=1pt,    
  scaled ticks=false,
  scaled y ticks=false,
  unbounded coords=jump,
  samples=101,
  minor tick length=0,
  scatter/classes={
    open={mark=*,draw=black,fill=white, mark size=2, thin},
    closed={mark=*,draw=black,fill=black, mark size=2, thin}
  },
  width=3in,
}
\pgfplotsset{open/.style={only marks, mark=*, fill=white, thin, mark size=2}}
\pgfplotsset{closed/.style={only marks, mark=*, draw=black, fill=black,
    thin, mark size=2}}
\pgfplotsset{labeled/.style={nodes near coords, only marks, mark=*, draw=black, fill=black, thin, mark size=2, point meta=explicit symbolic}}

\pgfplotsset{gridgraph/.style={clip=false, axis line style={shorten
      >=-8pt}, xlabel style={xshift=8pt}, ylabel style={yshift=8pt},
    enlargelimits=false, grid=major, tick style={black}}} 

\usepackage{tikz}
\usetikzlibrary{
  matrix,%
  % enables the snake paths
  decorations.pathmorphing,%
  % enables bigger arrows
  decorations.markings,%
  decorations.pathreplacing,%
  decorations.text,%
  calc,%
  patterns,%
  shapes.geometric,%
  arrows,
  decorations.shapes,
  positioning,
  plotmarks,
  intersections
}
\tikzset{>=latex} % sets default arrow type in tikz
\tikzset{font=\normalsize}
\tikzset{mark size=1.5pt, mark options=thin}
\tikzset{pin distance=4pt,
  every pin edge/.style={<-, thin, shorten <= -2pt}}

\interfootnotelinepenalty=10000
\renewcommand{\thefootnote}{(\arabic{footnote})}

\usepackage{multirow, bigdelim}

% change to \usepackage[sols]{handout} to see solutions
\usepackage{handout}

\newcommand\iscurrentenumi[1]{%
  \ifthenelse{\equal{#1}{\theenumi}}{}{#1}}
\setenumerate[1]{ref=\#\arabic*}
\setenumerate[2]{ref=\protect\iscurrentenumi{\theenumi}(\alph*)}
\newcommand\enumiiprefix[2]{%
  \ifthenelse{{\equal{#1}{\theenumi}}\AND{\equal{#2}{(\alph{enumii})}}}{}{}%
  \ifthenelse{{\equal{#1}{\theenumi}}\AND{\NOT{\equal{#2}{(\alph{enumii})}}}}{#2}{}%
  \ifthenelse{\equal{#1}{\theenumi}}{}{#1#2}%
}
\setenumerate[3]{ref=\protect\enumiiprefix{\theenumi}{(\alph{enumii})}\roman*}
% Make an environment that puts items in columns, first going across. 
\usepackage{tabto}
\newenvironment{enumeratecols}[2][]
{\NumTabs{#2}%
  \begin{enumerate*}[%before={\hspace{\dimexpr-\parindent-1pt}\tab},%
    itemjoin={\tab},#1]}%
  {\end{enumerate*}}

\newcommand{\related}[1]{%
\fcolorbox{black}{yellow}{%
  %\parbox[t]{12pt}{$\clubsuit$}%
  \parbox[t]{\linewidth-8pt}{\emph{Related problems:} #1}}}

\renewcommand{\answer}[1]{#1}

\definecolor{purple}{rgb}{0.5,0,1.0}

\usepackage{amsthm}

% make URLs print in the same font
\usepackage{url}
\urlstyle{same}

\usepackage{pifont}
\def\exend{\hfill\ding{118}}

%\makeatletter\@addtoreset{equation}{section}\makeatother
%\renewcommand{\theequation}{\arabic{section}.\arabic{equation}}

\newtheoremstyle{theorem}% name
{8pt}%      Space above
{0pt}%      Space below
{\itshape}%         Body font
{}%         Indent amount (empty = no indent, \parindent = para indent)
{\bfseries}% Thm head font
{.}%        Punctuation after thm head
{.5em}%     Space after thm head: " " = normal interword space;
                                %       \newline = linebreak
{}%         Thm head spec (can be left empty, meaning `normal')

\theoremstyle{theorem}

\let\Theorem\undefined
\newtheorem{Theorem}[equation]{Theorem}
\let\Definition\undefined
\newtheorem{Definition}[equation]{Definition}
\let\Summary\undefined
\newtheorem{Summary}[equation]{Summary}
\let\Conjecture\undefined
\newtheorem{Conjecture}[equation]{Conjecture}
\let\Fact\undefined
\newtheorem{Fact}[equation]{Fact}

\renewcommand{\thesubsection}{\arabic{subsection}}
\makeatletter
\def\@seccntformat#1{\@ifundefined{#1@cntformat}%
   {\csname the#1\endcsname\quad}%       default
   {\csname #1@cntformat\endcsname}}%    enable individual control
\newcommand\section@cntformat{}
\makeatother

  \usepackage{exam}

\usepackage{fancyhdr}
\lhead{\textcolor{white}{This content is protected and may not be shared, uploaded, or distributed.}}
\rhead{}
\rfoot{\vspace*{\fill}\footnotesize\textcopyright{} \emph{Harvard Math Ma}}
\renewcommand{\headrulewidth}{0pt}
\pagestyle{fancy}
\begin{document}

  \title{Practice Mini-Exam 3}

\instructions{instructions.tex}{}

\listofproblems

\begin{enumerate}
\item \points{12} \textbf{Algebra}.

  \begin{enumerate}
  \item 
    \begin{problem}[1in]
      Simplify $25^{-3/2}$. 
    \end{problem}

    \begin{solution}
      $25^{-3/2} = \dfrac{1}{25^{3/2}} = \dfrac{1}{(25^{1/2})^3} = \dfrac{1}{5^3} = \boxed{\frac{1}{125}}$ 
    \end{solution}

  \item 
    \begin{problem}[\stretch{1.5}]
      Simplify $\dfrac{\dfrac{1}{(a + b)^2} - \dfrac{1}{a^2}}{b}$ (assuming the denominators $a, b, a + b$ are non-zero).
    \end{problem}

    \begin{solution}
      First, let's simplify the numerator. The first step is to get everything over a common denominator. 
      \begin{align*}
        \dfrac{1}{(a + b)^2} - \dfrac{1}{a^2} &= \dfrac{a^2}{ a^2 (a + b)^2} - \dfrac{(a + b)^2}{a^2 (a + b)^2} \\
        &= \dfrac{a^2 - (a + b)^2}{a^2 (a + b)^2} \\
        &= \dfrac{a^2 - (a^2 + 2ab + b^2)}{a^2 (a + b)^2} \\
        &= \dfrac{-2ab-b^2}{a^2 (a + b)^2}         
      \end{align*}
      Now that we've simplified the numerator, we can go back to the original expression. 
      \begin{align*}
        \dfrac{\dfrac{1}{(a + b)^2} - \dfrac{1}{a^2}}{b} &= \dfrac{\quad \dfrac{-2ab-b^2}{a^2 (a + b)^2}\quad}{b} \\
        &= \dfrac{-2ab-b^2}{a^2 (a + b)^2} \cdot \frac{1}{b} \\
        &= \dfrac{b (-2a-b)}{a^2 (a + b)^2} \cdot \frac{1}{b} \\
        &= \boxed{\dfrac{-2a-b}{a^2 (a + b)^2}}
      \end{align*}
    \end{solution}

  \item \label{fall2017-practice-gateway-3-13}     

    \begin{problem}[\vfill\newpage]
      Solve $\dfrac{M}{a + Mb} = 2ab$ for $M$. 
    \end{problem}

    \begin{solution}
      Let's first multiply both sides by $a + Mb$ to get rid of the fraction: 
      \begin{align*}
        M & = (a + Mb)(2ab) \\
        \intertext{This is linear in $M$ (if you aren't sure what this means, see \href{https://people.math.harvard.edu/~jjchen/mathma/handouts/solving-equations.html}{the ``Algebra Practice: Solving equations'' handout}), so we'll try to rewrite it as $(\text{something}) M = \text{something}$.}
        M &= 2a^2 b + 2a b^2 M \\
        M - 2ab^2 M &= 2a^2 b \\
        (1 - 2ab^2) M & = 2a^2b \\
        M & = \boxed{\frac{2a^2b}{1 - 2ab^2}}
      \end{align*}
    \end{solution}

  \item

    \begin{problem}[\vfill\newpage]
      Solve $\dfrac{\dfrac{3}{x} - 4x}{\dfrac{1}{x} + 1} = 2$. 
    \end{problem}

    \begin{solution}
      Let's first multiply both sides by $\dfrac{1}{x} + 1$ to get rid of one denominator:
      \begin{align*}
        \dfrac{3}{x} - 4x &= \dfrac{2}{x} + 2 \\
        \dfrac{1}{x} - 4x &= 2 \\
        \intertext{Let's multiply both sides by $x$ to get rid of the remaining denominator:}
        1 - 4x^2 &= 2x \\
        \intertext{This isn't linear in $x$, so we'll move everything to one side and try to factor:}
        0 &= 4x^2 + 2x - 1 
      \end{align*}
      I don't see how to factor this, so I'll use the quadratic formula instead: \[x = \dfrac{-2 \pm \sqrt{2^2 - 4 (4)(-1)}}{8} = \dfrac{-2 \pm \sqrt{20}}{8} = \dfrac{-2 \pm 2 \sqrt{5}}{8} = \boxed{\frac{-1 \pm \sqrt{5}}{4}}.\]
    \end{solution}
  \end{enumerate}

\item \label{fall2018-midterm-5} \points{6} %CJ

  \def\FunctionF(#1){0.1*((#1)^3- 9*(#1))}%
  \def\FunctionG(#1){0.2*((#1)^3- 9*(#1))}%
  \def\FunctionH(#1){0.1*(((#1))^3- 9*((#1)))-2}%
  \def\FunctionP(#1){0.1*((2*(#1))^3- 9*(2*(#1)))}%
  \def\FunctionQ(#1){0.1*((2*(#1))^3- 9*(2*(#1)))-2}%
  \def\FunctionR(#1){-1*(0.1*((#1)^3- 9*(#1)))}%

  Here's the graph of a function $f(x)$. 

  \begin{center}	
    \begin{tikzpicture}
      \begin{axis}[
        xmin=-5.5, 
        xmax=5.5, 
        ymin=-5.5, 
        ymax=5.5, 
        grid=both, 
        xtick={-4, -2, ..., 4},
        minor xtick={-5, -4, ..., 5}, 
        ytick={-4, -2, ..., 4},
        minor ytick={-5, ..., 5},
        width=3in, axis equal image
        ]
        \addplot [domain=-5:5] {\FunctionF(x)};
      \end{axis}
    \end{tikzpicture}
  \end{center}

  Each picture below shows $f$ together with the graph of another function (the dashed one); the last graph is identical to $f$. Match each dashed graph with the appropriate function (choices A - H), and write your answer in the blank below the graph. No explanation is necessary. 

  \answeronly{Going across:}

  \begin{center}
    \begin{tabular}{ccccc}
      \begin{tikzpicture}
        \begin{axis}[width=2.25in, xmin=-5.5, xmax=5.5, ymin=-5.5, ymax=5.5, grid=both, xtick={-4, -2, 2, 4}, ytick={-4, -2, 2, 4}, minor xtick={-5, ..., 5}, minor ytick={-5, ..., 5}, ]
          \addplot+[domain=-5:5] {\FunctionF(x)}; %
          \addplot [domain=-5:5, dashed] {\FunctionG(x)};
        \end{axis}
      \end{tikzpicture} & \hspace{0.3cm} & 
      \begin{tikzpicture}
        \begin{axis}[width=2.25in, xmin=-5.5, xmax=5.5, ymin=-5.5, ymax=5.5, grid=both, xtick={-4, -2, 2, 4}, ytick={-4, -2, 2, 4}, minor xtick={-5, ..., 5}, minor ytick={-5, ..., 5}, ]
          \addplot+[domain=-5:5] {\FunctionF(x)}; %
          \addplot [domain=-5:5, dashed] {\FunctionH(x)};
        \end{axis}
      \end{tikzpicture} & \hspace{0.3cm} &
      \begin{tikzpicture}
        \begin{axis}[width=2.25in, xmin=-5.5, xmax=5.5, ymin=-5.5, ymax=5.5, grid=both, xtick={-4, -2, 2, 4}, ytick={-4, -2, 2, 4}, minor xtick={-5, ..., 5}, minor ytick={-5, ..., 5}, ]
          \addplot+[domain=-5:5] {\FunctionF(x)}; %
          \addplot [domain=-3:3, dashed] {\FunctionP(x)};
        \end{axis}
      \end{tikzpicture} \\[4pt]
      \cfitb{0.5in}{G} && \cfitb{0.5in}{A} && \cfitb{0.5in}{D} \\[24pt]
      \begin{tikzpicture}
        \begin{axis}[width=2.25in, xmin=-5.5, xmax=5.5, ymin=-5.5, ymax=5.5, grid=both, xtick={-4, -2, 2, 4}, ytick={-4, -2, 2, 4}, minor xtick={-5, ..., 5}, minor ytick={-5, ..., 5}, ]
          \addplot+[domain=-5:5] {\FunctionF(x)}; %
          \addplot [domain=-5:5, dashed] {\FunctionR(x)};
        \end{axis}
      \end{tikzpicture} &&
      \begin{tikzpicture}
        \begin{axis}[width=2.25in, xmin=-5.5, xmax=5.5, ymin=-5.5, ymax=5.5, grid=both, xtick={-4, -2, 2, 4}, ytick={-4, -2, 2, 4}, minor xtick={-5, ..., 5}, minor ytick={-5, ..., 5}, ]
          \addplot+[domain=-5:5] {\FunctionF(x)}; %
          \addplot [domain=-3:3, dashed] {\FunctionQ(x)};
        \end{axis}
      \end{tikzpicture} &&
      \begin{tikzpicture}
        \begin{axis}[width=2.25in, xmin=-5.5, xmax=5.5, ymin=-5.5, ymax=5.5, grid=both, xtick={-4, -2, 2, 4}, ytick={-4, -2, 2, 4}, minor xtick={-5, ..., 5}, minor ytick={-5, ..., 5}, ]
          \addplot [domain=-5:5, dashed] {\FunctionF(x)};
        \end{axis}
      \end{tikzpicture} \\[4pt]
      \cfitb{0.5in}{C} && \cfitb{0.5in}{E} && \cfitb{0.5in}{F} 
    \end{tabular}
  \end{center}

  Choices:

  \begin{enumeratecols}[label=\Alph*.]{4}
  \item $y=f(x)-2$
  \item $y=f\left(\frac{x}{2}\right)$
  \item $y=f(-x)$
  \item $y=f(2x)$
  \item $y=f(2x)-2$
  \item $y=-f(-x)$ 
  \item $y=2f(x)$ 
  \item $y=2f(x)-2$    
  \end{enumeratecols}

  \pspace{\newpage}

\item \label{fall2013-1a-midterm1-8} \points{6} % Omar

  Max leaves home at 8am for a job interview that will be in a town $60$ miles away. Traffic is bad at first, so he covers the first 30 miles at an average speed of $45$ miles per hour. Traffic improves after that, so he covers the remaining 30 miles at an average speed of $60$ miles per hour.

  \solonly{\related{\href{https://people.math.harvard.edu/~jjchen/mathma/docs/homework/PS03.html}{Problem Set 3}, {{\#5}}}}

  \begin{enumerate}
  \item
    \begin{problem}[\vfill]
      What is his average speed for the whole trip? Please include units in your answer. 
    \end{problem}

    \begin{solution}
      The average speed for Max's whole trip is the total distance ($60$ miles) divided by his total time. So, we need to figure out the total time his trip took. In the first part of the trip, he traveled 30 miles at an average speed of 45 miles per hour, taking $\dfrac{30 \text{ miles}}{45 \text{ miles per hour}} = \dfrac{2}{3}$ hours. In the second part of his trip, he traveled 30 miles at an average speed of 60 miles per hour, taking $\dfrac{30 \text{ miles}}{60 \text{ miles per hour}} = \dfrac{1}{2}$ of an hour. So, his 60 mile trip took a total of $\displaystyle \frac{2}{3}+\frac{1}{2} = \frac{7}{6}\ \text{hours}$. This means his average speed for the entire trip was $\dfrac{60\ \text{miles}}{7/6\ \text{hours}} = \boxed{\frac{360}{7} \text{ miles per hour}}$.
    \end{solution}

  \item
    \begin{problem}[\vfill\newpage]
      By speeding up even more on the second half of the drive, could Max have achieved an average speed of 100 miles per hour for the whole trip? If you think he could, what would his average speed for the second $30$ miles have to be? If you think he couldn't, explain carefully why not.
    \end{problem}

    \begin{solution}
      Again, let's think about the total time for his trip. In order to average 100 miles per hour for the entire 60 mile trip, he would have had to travel 60 miles in $\dfrac{60 \text{ miles}}{100 \text{ miles per hour}} = 0.6$ hours. However, we already calculated that he spent $\frac{2}{3}$ hours traveling the first 30 miles, so he certainly \fbox{can't} travel the entire 60 miles in less than that!
    \end{solution}

  \end{enumerate}

\item \points{6} 

  North Dakota's Highway $46$ is supposedly the straightest road in the United States. From the city of Streeter, it runs $123$ miles east to Hickson. Litchville is a tiny city on this highway.

  A bicyclist decides to ride on the highway one day. Let $s(t)$ be her position at time $t$, where $t$ is measured in hours after noon and position is given in miles east of Litchville. Here's the graph of $s(t)$:

  \begin{center}
    \begin{tikzpicture}[declare function={f(\x)=3-(\x-1)*(\x-3)*(\x-4)*(\x-7)/5;}]
      \begin{axis}[gridgraph, xlabel=$t$, ylabel=position, width=4.5in, 
        height=2.5in, xtick={1, 2, ..., 8}, ytick={-16, -12, ..., 12}, ymin=-16, ymax=12, xmax=8]
        \addplot+[domain=0:7.5] {f(x)};
        \addplot+[closed] coordinates {(7.5, {f(7.5)})}; 
      \end{axis}
    \end{tikzpicture}
  \end{center}

  \begin{enumerate}
  \item
    \begin{problem}[\vfill]
      At roughly what times is the cyclist heading west? 
    \end{problem}

    \begin{solution}
      We're told that $s(t)$ gives the cyclist's position in miles east of Litchville. That means that, when $s(t)$ is increasing, the cyclist is going east; when $s(t)$ is decreasing, the cyclist is going west.

      It looks like $s(t)$ is decreasing on roughly $(1.7, 3.5)$ and $(6, 7.5)$. So, the times when the cyclist is going west are \fbox{around 1:45 pm - 3:30 pm and 6 pm - 7:30 pm}. 
    \end{solution}

  \item
    \begin{problem}[\vfill]
      Your friend is looking at the cyclist's average velocity between 5 pm and 7 pm. She remembers that there's some way to visualize this on the graph of $s(t)$, but she doesn't remember exactly what it is. Explain it to her, and illustrate on the graph above. 
    \end{problem}

    \begin{solution}
      The cyclist's average velocity between 5 pm and 7 pm is $\dfrac{s(7) - s(5)}{7 - 5}$, which we can picture as the slope of this secant line:
      \begin{center}
        \begin{tikzpicture}[declare function={f(\x)=3-(\x-1)*(\x-3)*(\x-4)*(\x-7)/5;}]
          \begin{axis}[gridgraph, xlabel=$t$, ylabel=position, width=4.5in, 
            height=2.5in, xtick={1, 2, ..., 8}, ytick={-16, -12, ..., 12}, ymin=-16, ymax=12, xmax=8]
            \addplot+[domain=0:7.5] {f(x)}; %
            \addplot+[closed] coordinates {(7.5, {f(7.5)})}; %
            \addplot[sharp plot, red, samples=2] coordinates { (5, {f(5)}) (7, {f(7)})}; % 
          \end{axis}
        \end{tikzpicture}
      \end{center}      
    \end{solution}

  \item
    \begin{problem}[\vfill\newpage]
      When was the cyclist furthest from Litchfield? 
    \end{problem}

    \begin{solution}
      We want to know when the cyclist was furthest east or west of Litchfield. From the graph, she was furthest west at the start of her ride, when $s$ is about $-14$ (so she was about 14 miles west). She was furthest east roughly at $t = 6$, when she was about 9 miles east of Litchfield. So, the cyclist was furthest from Litchfield \fbox{at noon}. 
    \end{solution}

  \end{enumerate}

\item \points{6}  \finishexam

  \begin{problem}
    Sketch the graph of $\displaystyle f(x) = \left| \dfrac{1}{x + 1} - 2 \right|$, label any intercepts and asymptotes with their exact coordinates. (Please show your work to indicate how you found them.)

    We encourage you to sketch some intermediate steps so you can get partial credit if your final sketch is incorrect.
  \end{problem}

  \begin{solution}
    We can start by graphing $\dfrac{1}{x + 1} - 2$, which we get by starting with $y = \dfrac{1}{x}$ and shifting it left by 1 and down by 2, so its asymptotes are now at $x = -1$ and $y = -2$:
    \begin{center}
      \begin{tikzpicture}
        \begin{axis}[ymin=-5, ymax=5, xtick={-1}, ytick={-2}]
          \addplot+[domain=-5:5] {(1/(x+1)-2)}; %
          \addplot+[domain=-5:5, dashed] {-2}; %
          \addplot+[domain=-5:5, dashed] (-1, x); % 
        \end{axis}
      \end{tikzpicture}
    \end{center}
    Now, we take the absolute value. When $\dfrac{1}{x + 1} - 2$ is $\geq 0$, the absolute value is just the same as the original. When $\dfrac{1}{x + 1} - 2$ is negative, the absolute value negates it, which flips those parts of the graph over the $x$-axis. This doesn't change the vertical asymptote but moves the horizontal asymptote to $y = 2$:
    \begin{center}
      \begin{tikzpicture}
        \begin{axis}[ymin=-5, ymax=5, xtick={-1}, ytick={2}]
          \addplot+[domain=-5:5] {abs(1/(x+1)-2)}; 
          \addplot+[domain=-5:5, dashed] {2}; %
          \addplot+[domain=-5:5, dashed] (-1, x); % 
        \end{axis}
      \end{tikzpicture}
    \end{center}
    The $y$-intercept is $f(0) = \left| \dfrac{1}{1} - 2 \right| = 1$. The $x$-intercepts of $f$ are the same as the $x$-intercepts of $\dfrac{1}{x + 1} - 2$, and we can solve for these:
    \begin{align*}
      \dfrac{1}{x + 1} - 2 &= 0 \\
      \dfrac{1}{x + 1} &= 2 \\
      \dfrac{1}{2} &= x + 1 \\
      - \dfrac{1}{2} &= x 
    \end{align*}
    So, here's our labeled graph:
    \begin{center}
      \answer{
        \begin{tikzpicture}
          \begin{axis}[ymin=-5, ymax=5, xtick={-1, -0.5}, xticklabels={$-1$, $-\frac{1}{2}$}, ytick={2, 1}, width=4in]
            \addplot+[domain=-5:5] {abs(1/(x+1)-2)}; 
            \addplot+[domain=-5:5, dashed] {2}; %
            \addplot+[domain=-5:5, dashed] (-1, x); %
          \end{axis}
        \end{tikzpicture}
      }
    \end{center}
  \end{solution}

\end{enumerate}

\end{document}
